\section*{Setup}

Every number in this report comes from running the code in this repository on
my laptop, and the handout asks for three details of it:

\begin{center}
\begin{tabular}{ll}
\toprule
CPU & Apple M4 Pro, 12 cores (8 performance + 4 efficiency)\\
RAM & 24 GiB\\
NumPy & 2.4.6\\
\bottomrule
\end{tabular}
\end{center}

\textbf{From scratch.} Everything a problem asks me to implement is written with
NumPy array arithmetic: the DFT matrix, all transfer functions and the filtering
pipeline, the homomorphic filter, the bilateral and Kuwahara filters, the Wiener
filters and the noise estimators, the motion PSF, the inverse and Wiener
deconvolution, and the blur estimator. I use \code{numpy.fft} for the transforms
where the handout allows it. The metrics are my own too: PSNR, SSIM, and the
connected-component counts in 2.3. SciPy appears only where the handout asks for
a comparison (\code{scipy.signal.convolve2d} and \code{fftconvolve} in 1.3,
\code{scipy.ndimage.gaussian\_filter} in 2.2(d)) and in \code{scipy.special.sici},
which only draws the predicted $\Si$ curve in Fig.~\ref{fig:ringing}.
scikit-image and \code{scipy.ndimage.label} are only used in the unit tests, to
check my SSIM and my component labelling. Pillow reads the images and Matplotlib
draws the figures.

\textbf{Inputs.} The supplied images are used as given: \code{q1.png}
(Problem 1), the four Problem 2 images, \code{tunnel\_image.jpg} (Problem 3),
\code{q4.png} (Problem 4) and \code{reference\_image.jpg} (Problems 5 and 6).
For 5.2 and 6.4 I took my own photos with a phone. Any preprocessing is stated
where it is applied. Random noise always comes from
\code{numpy.random.default\_rng} with the seed given in that section.

\textbf{Scope.} The marks table on page 1 of the handout lists a Problem 7 (lens
distortion), but the handout has no text for it, so I did not attempt it.

\textbf{Reproducing.} From the repository root:
\begin{itemize}
\item \code{python src/run\_experiments.py} runs every problem,
\item \code{python src/run\_experiments.py -{}-problem 2} runs one problem,
\item \code{python -m unittest discover -s tests -v} runs the 56 unit tests.
\end{itemize}
Each run writes its tables and figures to \code{results/problem\_k/} and its
terminal output to \code{results/logs/}. The \code{*\_terminal.png} screenshots
next to the results show those runs.
